% Unit: Computer Science A --- testing, mocks, and logging (\input after \texttt{cs-a/03-practice.tex}; lambdas in \texttt{cs-a/lambda.tex}, REST/Javalin in \texttt{cs-a/rest.tex}).

\runningheader{Computer Science A}{Unit 4: Testing, mocks, and logging}{Page \thepage\ of \numpages}

\begingradingrange{csau4}

\uplevel{\par\textbf{Testing and mocks}\par}

\question[4]
What is Mockito and why would we use it?
\makeemptybox{1.25in}
\begin{solution}
  Popular Java mocking library: stub dependencies, verify interactions, isolate the class under test without slow or brittle real I/O or databases.
\end{solution}

\question[4]
What is a mock?
\makeemptybox{1.05in}
\begin{solution}
  A test double with programmed behavior and expectations: returns canned values, records calls, and fails tests if expected interactions do not occur.
\end{solution}

\question[4]
What is test driven development?
\makeemptybox{1.2in}
\begin{solution}
  Write a failing test first, implement minimal code to pass, refactor; repeats in short cycles to drive design and regression safety.
\end{solution}

\question[4]
Why are unit tests important?
\makeemptybox{1.25in}
\begin{solution}
  Catch regressions early, document expected behavior, enable safe refactoring, shorten debug loops, and support continuous integration.
\end{solution}

\question[4]
How can JUnit annotations help with running our tests?
\makeemptybox{1.35in}
\begin{solution}
  Mark lifecycle (\texttt{@BeforeEach}, \texttt{@AfterEach}), identify tests (\texttt{@Test}), disable/skip, parameterize, order, tag for suites, and set timeouts or expected exceptions.
\end{solution}

\question[4]
Give me an example of how you would test a method that takes in two integers as input and returns an integer representing the sum.
\makeemptybox{1.55in}
\begin{solution}
  Call the method with known pairs and assert the result, e.g.\ \texttt{assertEquals(5, add(2, 3));} plus edge cases (\texttt{0}, negatives) using JUnit \texttt{assertions}.
\end{solution}

\newpage

\uplevel{\par\textbf{Logging}\par}

\question[4]
What is logging and what are the benefits of it?
\makeemptybox{1.25in}
\begin{solution}
  Recording runtime events (errors, timing, decisions) to appenders (console, files, aggregators). Benefits: production diagnosis without debugger, audit trails, and structured observability.
\end{solution}

\question[4]
Describe the different logging levels and how they should be used.
\makeemptybox{1.45in}
\begin{solution}
  Typical levels (e.g.\ SLF4J): \texttt{ERROR}/\texttt{WARN} for problems, \texttt{INFO} for high-signal lifecycle events, \texttt{DEBUG}/\texttt{TRACE} for developer detail---tune thresholds per environment to limit noise and cost.
\end{solution}

\question[4]
How would we configure logging in an application?
\makeemptybox{1.35in}
\begin{solution}
  Use a facade (\texttt{slf4j}) with a binding (Logback, Log4j2): \texttt{logback.xml} or \texttt{log4j2.xml} for appenders, levels per logger, patterns, and optional env-specific profiles.
\end{solution}

\endgradingrange{csau4}
